\chapter{2020年全国III卷（文）}

\section{单选题}

\begin{problem}
已知集合 \(A=\{1,2,3,5,7,11\}\)，\(B=\{x\mid 3<x<15\}\)，则 \(A\cap B\) 中元素的个数为（\quad）

\choices
    {\(2\)}
    {\(3\)}
    {\(4\)}
    {\(5\)}

\end{problem}

\begin{answer}
B
\end{answer}

\begin{solution}
由集合 \(A=\{1,2,3,5,7,11\}\)，\(B=\{x\mid 3<x<15\}\)，得 \(A\cap B=\{5,7,11\}\)，所以 \(A\cap B\) 中元素的个数为 \(3\). 故选 B.
\end{solution}
\begin{problem}
若 \(\overline{z}(1+\i)=1-\i\)，则 \(z=\)（\quad）

\choices
    {\(1-\i\)}
    {\(1+\i\)}
    {\(-\i\)}
    {\(\i\)}

\end{problem}

\begin{answer}
D
\end{answer}

\begin{solution}
由 \(\overline{z}(1+\i)=1-\i\)，得
\(\overline{z}=\frac{1-\i}{1+\i}=\frac{(1-\i)^2}{(1+\i)(1-\i)}=-\i\).

所以 \(z=\i\). 故选 D.
\end{solution}
\begin{problem}
设一组样本数据 \(x_1,x_2,\cdots,x_n\) 的方差为 \(0.01\)，则数据 \(10x_1,10x_2,\cdots,10x_n\) 的方差为（\quad）

\choices
    {\(0.01\)}
    {\(0.1\)}
    {\(1\)}
    {\(10\)}

\end{problem}

\begin{answer}
C
\end{answer}

\begin{solution}
样本数据 \(x_1,x_2,\cdots,x_n\) 的方差为 \(0.01\). 根据任何一组数据同时扩大几倍，方差将变为平方倍增长，数据 \(10x_1,10x_2,\cdots,10x_n\) 的方差为 \(100\times 0.01=1\). 故选 C.
\end{solution}
\begin{problem}
Logistic 模型是常用数学模型之一，可应用于流行病学领域. 有学者根据公布数据建立了某地区新冠肺炎累计确诊病例数 \(I(t)\)（\(t\) 的单位：天）的 Logistic 模型：\(I(t)=\dfrac{K}{1+\e^{-0.23(t-53)}}\)，其中 \(K\) 为最大确诊病例数. 当 \(I(t^*)=0.95K\) 时，标志着已初步遏制疫情，则 \(t^*\) 约为（\quad）（\(\ln 19\approx 3\)）

\choices
    {\(60\)}
    {\(63\)}
    {\(66\)}
    {\(69\)}

\end{problem}

\begin{answer}
C
\end{answer}

\begin{solution}
由已知可得
\(\frac{K}{1+\e^{-0.23(t^*-53)}}=0.95K\)，
解得
\(\e^{-0.23(t^*-53)}=\frac{1}{19}\).

两边取对数有
\(-0.23(t^*-53)=-\ln 19\).

所以 \(t^*\approx 66\). 故选 C.
\end{solution}
\begin{problem}
已知 \(\sin\theta+\sin\left( \theta+\dfrac{\pi}{3} \right)=1\)，则 \(\sin\left( \theta+\dfrac{\pi}{6} \right)=\)（\quad）

\choices
    {\(\dfrac{1}{2}\)}
    {\(\dfrac{\sqrt3}{3}\)}
    {\(\dfrac{2}{3}\)}
    {\(\dfrac{\sqrt2}{2}\)}

\end{problem}

\begin{answer}
B
\end{answer}

\begin{solution}
由
\(\sin\theta+\sin\left( \theta+\frac{\pi}{3} \right)=1\)，
得
\(\sin\theta+\frac{1}{2}\sin\theta+\frac{\sqrt3}{2}\cos\theta=1\)，
即
\(\frac{3}{2}\sin\theta+\frac{\sqrt3}{2}\cos\theta=1\).

所以
\(\sqrt3\sin\left( \theta+\frac{\pi}{6} \right)=1\)，
从而
\(\sin\left( \theta+\frac{\pi}{6} \right)=\frac{\sqrt3}{3}\).

故选 B.
\end{solution}
\begin{problem}
在平面内，\(A,B\) 是两个定点，\(C\) 是动点. 若 \(\overrightarrow{AC}\cdot\overrightarrow{BC}=1\)，则点 \(C\) 的轨迹为（\quad）

\choices
    {圆}
    {椭圆}
    {抛物线}
    {直线}

\end{problem}

\begin{answer}
A
\end{answer}

\begin{solution}
不妨设 \(A(-a,0)\)，\(B(a,0)\)，设 \(C(x,y)\). 因为
\(\overrightarrow{AC}\cdot\overrightarrow{BC}=1\)，
所以
\((x+a,y)\cdot(x-a,y)=1\).

解得
\(x^2+y^2=a^2+1\).

所以点 \(C\) 的轨迹为圆. 故选 A.
\end{solution}
\begin{problem}
设 \(O\) 为坐标原点，直线 \(x=2\) 与抛物线 \(C:y^2=2px\)（\(p>0\)）交于 \(D,E\) 两点，若 \(OD\perp OE\)，则 \(C\) 的焦点坐标为（\quad）

\choices
    {\(\left( \dfrac{1}{4},0 \right)\)}
    {\(\left( \dfrac{1}{2},0 \right)\)}
    {\((1,0)\)}
    {\((2,0)\)}

\end{problem}

\begin{answer}
B
\end{answer}

\begin{solution}
将 \(x=2\) 代入抛物线 \(y^2=2px\)，可得 \(y=\pm 2\sqrt{p}\). 因为 \(OD\perp OE\)，可得 \(k_{OD}\cdot k_{OE}=-1\)，即
\(\frac{2\sqrt{p}}{2}\cdot\frac{-2\sqrt{p}}{2}=-1\)，
解得 \(p=1\).

所以抛物线方程为 \(y^2=2x\)，它的焦点坐标为 \(\left( \frac{1}{2},0 \right)\). 故选 B.
\end{solution}
\begin{problem}
点 \((0,-1)\) 到直线 \(y=k(x+1)\) 距离的最大值为（\quad）

\choices
    {\(1\)}
    {\(\sqrt2\)}
    {\(\sqrt3\)}
    {\(2\)}

\end{problem}

\begin{answer}
B
\end{answer}

\begin{solution}
点 \((0,-1)\) 到直线 \(y=k(x+1)\) 的距离
\(d=\frac{|k+1|}{\sqrt{k^2+1}}=\sqrt{\frac{k^2+2k+1}{k^2+1}}=\sqrt{1+\frac{2k}{k^2+1}}\).

要求距离的最大值，故需 \(k>0\). 由
\(k^2+1\ge 2k\)
可得
\(d\le \sqrt{1+\frac{2k}{2k}}=\sqrt2\).

当 \(k=1\) 时等号成立，所以距离的最大值为 \(\sqrt2\). 故选 B.
\end{solution}
\begin{problem}
如图为某几何体的三视图，则该几何体的表面积是（\quad）

\begin{center}
\bitmapinclude[max width=3.4cm,max totalheight=3cm]{img/2020/national_paper_3_liberal/q9_fig1.png}
\end{center}

\choices
    {\(6+4\sqrt2\)}
    {\(4+4\sqrt2\)}
    {\(6+2\sqrt3\)}
    {\(4+2\sqrt3\)}

\end{problem}

\begin{answer}
C
\end{answer}

\begin{solution}
由三视图可知，几何体的直观图是正方体的一个角. 设三棱锥为 \(P-ABC\)，则
\(PA=AB=AC=2\)，
且 \(PA,AB,AC\) 两两垂直.

故
\(PB=BC=PC=2\sqrt2\).

几何体的表面积为
\(3\times \frac{1}{2}\times 2\times 2+\frac{\sqrt3}{4}\times \left( 2\sqrt2 \right)^2=6+2\sqrt3\).

故选 C.
\end{solution}
\begin{problem}
设 \(a=\log_3 2\)，\(b=\log_5 3\)，\(c=\dfrac{2}{3}\)，则（\quad）

\choices
    {\(a<c<b\)}
    {\(a<b<c\)}
    {\(b<c<a\)}
    {\(c<a<b\)}

\end{problem}

\begin{answer}
A
\end{answer}

\begin{solution}
因为
\(a=\log_3 2=\log_3 \sqrt[3]{8}<\log_3 \sqrt[3]{9}=\frac{2}{3}\)，
且
\(b=\log_5 3=\log_5 \sqrt[3]{27}>\log_5 \sqrt[3]{25}=\frac{2}{3}\)，

所以
\(a<c<b\).

故选 A.
\end{solution}
\begin{problem}
在 \(\triangle ABC\) 中，\(\cos C=\dfrac{2}{3}\)，\(AC=4\)，\(BC=3\)，则 \(\tan B=\)（\quad）

\choices
    {\(\sqrt5\)}
    {\(2\sqrt5\)}
    {\(4\sqrt5\)}
    {\(8\sqrt5\)}

\end{problem}

\begin{answer}
C
\end{answer}

\begin{solution}
由 \(\cos C=\dfrac{2}{3}\)，得
\(\tan C=\sqrt{\frac{1}{\cos^2 C}-1}=\frac{\sqrt5}{2}\).

由余弦定理可得
\(AB=\sqrt{AC^2+BC^2-2AC\cdot BC\cdot \cos C} =\sqrt{4^2+3^2-2\times 4\times 3\times \frac{2}{3}}=3\)，
故 \(AB=BC\)，可得 \(\angle A=\angle C\).

所以
\(B=\pi-2C\)，
从而
\(\tan B=\tan(\pi-2C)=-\tan 2C=-\frac{2\tan C}{1-\tan^2 C}=4\sqrt5\).

故选 C.
\end{solution}
\begin{problem}
已知函数 \(f(x)=\sin x+\dfrac{1}{\sin x}\)，则（\quad）

\choices
    {\(f(x)\) 的最小值为 \(2\)}
    {\(f(x)\) 的图象关于 \(y\) 轴对称}
    {\(f(x)\) 的图象关于直线 \(x=\pi\) 对称}
    {\(f(x)\) 的图象关于直线 \(x=\dfrac{\pi}{2}\) 对称}

\end{problem}

\begin{answer}
D
\end{answer}

\begin{solution}
由 \(\sin x\neq 0\) 可得函数的定义域为 \(\{x\mid x\neq k\pi,\ k\in \Z\}\)，故定义域关于原点对称.

设 \(\sin x=t\)，则
\(y=f(x)=t+\frac{1}{t}\)，\(t\in[-1,1]\).
由双勾函数的图象和性质得，\(y\ge 2\) 或 \(y\le -2\)，故 A 错误.

又有
\(f(-x)=\sin(-x)+\frac{1}{\sin(-x)}=-\left( \sin x+\frac{1}{\sin x} \right)=-f(x)\)，
故 \(f(x)\) 是奇函数，且定义域关于原点对称，所以图象关于原点中心对称，B 错误.

由
\(f(\pi+x)=\sin(\pi+x)+\frac{1}{\sin(\pi+x)}=-\sin x-\frac{1}{\sin x}\)，
\(f(\pi-x)=\sin(\pi-x)+\frac{1}{\sin(\pi-x)}=\sin x+\frac{1}{\sin x}\)，
可知 \(f(\pi+x)\neq f(\pi-x)\)，故 \(f(x)\) 的图象不关于直线 \(x=\pi\) 对称，C 错误.

又 \(f\left( \frac{\pi}{2}+x \right)=\sin\left( \frac{\pi}{2}+x \right)+\frac{1}{\sin\left( \frac{\pi}{2}+x \right)}=\cos x+\frac{1}{\cos x}\)，
\(f\left( \frac{\pi}{2}-x \right)=\sin\left( \frac{\pi}{2}-x \right)+\frac{1}{\sin\left( \frac{\pi}{2}-x \right)}=\cos x+\frac{1}{\cos x}\)，
故
\(f\left( \frac{\pi}{2}+x \right)=f\left( \frac{\pi}{2}-x \right)\).

所以 \(f(x)\) 的图象关于直线 \(x=\dfrac{\pi}{2}\) 对称. 故选 D.
\end{solution}
\section{填空题}

\begin{problem}
若 \(x,y\) 满足约束条件 \(\left\{
\begin{array}{l}
x+y\ge 0,\\
2x-y\ge 0,\\
x\le 1,
\end{array}
\right.\) 则 \(z=3x+2y\) 的最大值为 \fillinblank{}.

\end{problem}

\begin{answer}
7
\end{answer}

\begin{solution}
\bitmapfigure[max width=4.6cm,max totalheight=4.2cm]{img/2020/national_paper_3_liberal/q13_solution_fig1.png}

先根据约束条件画出可行域，由
\examdisplaycases{}{
x=1,\\
2x-y=0
}
解得 \(A(1,2)\).

如图，当直线 \(z=3x+2y\) 过点 \(A(1,2)\) 时，目标函数在 \(y\) 轴上的截距取得最大值，此时 \(z\) 取得最大值.

所以当 \(x=1\)，\(y=2\) 时，
\(z_{\max}=3\times 1+2\times 2=7\).

故答案为：7.
\end{solution}
\begin{problem}
设双曲线 \(C:\dfrac{x^2}{a^2}-\dfrac{y^2}{b^2}=1\)（\(a>0\)，\(b>0\)）的一条渐近线为 \(y=\sqrt2x\)，则 \(C\) 的离心率为 \fillinblank{}.

\end{problem}

\begin{answer}
\(\sqrt3\)
\end{answer}

\begin{solution}
由双曲线的方程可得渐近线方程为
\(y=\pm \frac{b}{a}x\).

由题意可得
\(\frac{b}{a}=\sqrt2\)，
所以双曲线的离心率
\(e=\frac{c}{a}=\sqrt{1+\frac{b^2}{a^2}}=\sqrt3\).

故答案为：\(\sqrt3\).
\end{solution}
\begin{problem}
设函数 \(f(x)=\dfrac{\e^x}{x+a}\)，若 \(f'(1)=\dfrac{\e}{4}\)，则 \(a=\fillinblank{}\).

\end{problem}

\begin{answer}
1
\end{answer}

\begin{solution}
因为
\(f(x)=\frac{\e^x}{x+a}\)，
所以
\(f'(x)=\frac{(x+a-1)\e^x}{(x+a)^2}\).

由 \(f'(1)=\dfrac{\e}{4}\)，得
\(\frac{a\e}{(a+1)^2}=\frac{\e}{4}\)，
即
\(\frac{a}{(a+1)^2}=\frac{1}{4}\).

解得 \(a=1\). 故答案为：1.
\end{solution}
\begin{problem}
已知圆锥的底面半径为 \(1\)，母线长为 \(3\)，则该圆锥内半径最大的球的体积为 \fillinblank{}.

\end{problem}

\begin{answer}
\(\dfrac{\sqrt2\pi}{3}\)
\end{answer}

\begin{solution}
\bitmapfigure[max width=4.8cm,max totalheight=4.2cm]{img/2020/national_paper_3_liberal/q16_solution_fig1.png}

当球为该圆锥内切球时，半径最大.

如图，\(BS=3\)，\(BC=1\)，则圆锥高
\(SC=\sqrt{BS^2-BC^2}=\sqrt{9-1}=2\sqrt2\).

设内切球与圆锥相切于点 \(D\)，半径为 \(r\)，则 \(\triangle SOD\sim \triangle SCB\)，故有
\(\frac{SO}{BS}=\frac{OD}{BC}\)，
即
\(\frac{2\sqrt2-r}{3}=\frac{r}{1}\).

解得
\(r=\frac{\sqrt2}{2}\).

所以该球的体积为
\(\frac{4}{3}\pi r^3=\frac{4}{3}\pi\times \left( \frac{\sqrt2}{2} \right)^3=\frac{\sqrt2\pi}{3}\).

故答案为：\(\dfrac{\sqrt2\pi}{3}\).
\end{solution}
\section{解答题}

\begin{problem}
设等比数列 \(\{a_n\}\) 满足 \(a_1+a_2=4\)，\(a_3-a_1=8\).

\begin{enumerate}
    \item 求 \(\{a_n\}\) 的通项公式；
    \item 记 \(S_n\) 为数列 \(\{\log_3 a_n\}\) 的前 \(n\) 项和. 若 \(S_m+S_{m+1}=S_{m+3}\)，求 \(m\).
\end{enumerate}

\end{problem}

\begin{solution}
（1）设公比为 \(q\)，则由
\examdisplaycases{}{
a_1+a_1q=4,\\
a_1q^2-a_1=8
}
可得 \(a_1=1\)，\(q=3\).

所以
\(a_n=3^{n-1}\).

（2）由（1）有
\(\log_3 a_n=n-1\)，
是一个以 \(0\) 为首项、\(1\) 为公差的等差数列.

所以
\(S_n=\frac{n(n-1)}{2}\).

由已知可得
\(\frac{m(m-1)}{2}+\frac{(m+1)m}{2}=\frac{(m+3)(m+2)}{2}\)，
即
\(m^2-5m-6=0\).

解得 \(m=6\) 或 \(m=-1\)（舍去），所以 \(m=6\).
\end{solution}
\begin{problem}
某学生兴趣小组随机调查了某市 \(100\) 天中每天的空气质量等级和当天到某公园锻炼的人次，整理数据得到下表（单位：天）：

\newcommand{\airqualityheader}{%
\begin{tikzpicture}[baseline=(base.base)]
\node[inner sep=0pt,outer sep=0pt,minimum width=2.45cm,minimum height=1.05cm] (base) {};
\draw (base.north west) -- (base.south east);
\node[anchor=north east,inner sep=0pt,font=\zihao{-5}] at ([xshift=-1pt,yshift=-1pt]base.north east) {锻炼人次};
\node[anchor=south west,inner sep=0pt,font=\zihao{-5}] at ([xshift=1pt,yshift=1pt]base.south west) {空气质量等级};
\end{tikzpicture}%
}

\begin{center}
\begin{tabular}{c|ccc}
\airqualityheader & \([0,200]\) & \((200,400]\) & \((400,600]\) \\
\hline
\(1\)（优） & \(2\) & \(16\) & \(25\) \\
\(2\)（良） & \(5\) & \(10\) & \(12\) \\
\(3\)（轻度污染） & \(6\) & \(7\) & \(8\) \\
\(4\)（中度污染） & \(7\) & \(2\) & \(0\)
\end{tabular}
\end{center}

\begin{enumerate}
    \item 分别估计该市一天的空气质量等级为 \(1,2,3,4\) 的概率；
    \item 求一天中到该公园锻炼的平均人次的估计值（同一组中的数据用该组区间的中点值为代表）；
    \item 若某天的空气质量等级为 \(1\) 或 \(2\)，则称这天“空气质量好”；若某天的空气质量等级为 \(3\) 或 \(4\)，则称这天“空气质量不好”. 根据所给数据，完成下面的 \(2\times 2\) 列联表，并根据列联表，判断是否有 \(95\%\) 的把握认为一天中到该公园锻炼的人次与该市当天的空气质量有关？
\end{enumerate}

\begin{center}
\begin{tabular}{c|cc}
 & 人次\(\le 400\) & 人次\(>400\) \\
\hline
空气质量好 &  &  \\
空气质量不好 &  &
\end{tabular}
\end{center}

附：\(K^2=\dfrac{n(ad-bc)^2}{(a+b)(c+d)(a+c)(b+d)}\).

\begin{center}
\begin{tabular}{c|ccc}
\(P(K^2\ge k)\) & \(0.050\) & \(0.010\) & \(0.001\) \\
\hline
\(k\) & \(3.841\) & \(6.635\) & \(10.828\)
\end{tabular}
\end{center}

\end{problem}

\begin{solution}
（1）该市一天的空气质量等级为 \(1\) 的概率为
\(\frac{2+16+25}{100}=\frac{43}{100}\)，
该市一天的空气质量等级为 \(2\) 的概率为
\(\frac{5+10+12}{100}=\frac{27}{100}\)，
该市一天的空气质量等级为 \(3\) 的概率为
\(\frac{6+7+8}{100}=\frac{21}{100}\)，
该市一天的空气质量等级为 \(4\) 的概率为
\(\frac{7+2+0}{100}=\frac{9}{100}\).

（2）由题意可得，一天中到该公园锻炼的平均人次的估计值为
\(\bar{x}=100\times 0.20+300\times 0.35+500\times 0.45=350\).

（3）根据所给数据，可得下面的 \(2\times 2\) 列联表：

\begin{center}
\begin{tabular}{c|cc|c}
 & 人次\(\le 400\) & 人次\(>400\) & 总计 \\
\hline
空气质量好 & \(33\) & \(37\) & \(70\) \\
空气质量不好 & \(22\) & \(8\) & \(30\) \\
\hline
总计 & \(55\) & \(45\) & \(100\)
\end{tabular}
\end{center}

由表中数据可得
\(K^2=\frac{100\times (33\times 8-37\times 22)^2}{70\times 30\times 55\times 45}\approx 5.820>3.841\).

所以有 \(95\%\) 的把握认为一天中到该公园锻炼的人次与该市当天的空气质量有关.
\end{solution}
\begin{problem}
如图，在长方体 \(ABCD-A_1B_1C_1D_1\) 中，点 \(E,F\) 分别在棱 \(DD_1,BB_1\) 上，且 \(2DE=ED_1\)，\(BF=2FB_1\). 证明：

\begin{center}
\vspace{-0.5em}
\bitmapinclude[max width=3.2cm,max totalheight=4.2cm]{img/2020/national_paper_3_liberal/q19_fig1.png}
\vspace{-0.5em}
\end{center}

\begin{enumerate}
    \item 当 \(AB=BC\) 时，\(EF\perp AC\)；
    \item 点 \(C_1\) 在平面 \(AEF\) 内.
\end{enumerate}

\end{problem}

\begin{solution}
（1）因为 \(ABCD-A_1B_1C_1D_1\) 是长方体，所以 \(BB_1\perp\) 平面 \(ABCD\)，而 \(AC\subset\) 平面 \(ABCD\)，所以 \(AC\perp BB_1\).

因为 \(ABCD-A_1B_1C_1D_1\) 是长方体，且 \(AB=BC\)，所以 \(ABCD\) 是正方形，所以 \(AC\perp BD\).

又 \(BD\cap BB_1=B\)，所以 \(AC\perp\) 平面 \(BB_1D_1D\).

又因为点 \(E,F\) 分别在棱 \(DD_1,BB_1\) 上，所以 \(EF\subset\) 平面 \(BB_1D_1D\)，所以 \(EF\perp AC\).

（2）取 \(AA_1\) 上靠近 \(A_1\) 的三等分点 \(M\)，连接 \(D_1M\)、\(C_1F\)、\(MF\).

因为点 \(E\) 在 \(DD_1\) 上，且 \(2DE=ED_1\)，所以 \(ED_1\parallel AM\)，且 \(ED_1=AM\)，所以四边形 \(AED_1M\) 为平行四边形，所以 \(D_1M\parallel AE\)，且 \(D_1M=AE\).

又因为 \(F\) 在 \(BB_1\) 上，且 \(BF=2FB_1\)，所以 \(A_1M\parallel FB_1\)，且 \(A_1M=FB_1\)，所以 \(A_1B_1FM\) 为平行四边形.

所以 \(FM\parallel A_1B_1\)，且 \(FM=A_1B_1\)，即 \(FM\parallel C_1D_1\)，且 \(FM=C_1D_1\)，所以 \(C_1D_1MF\) 为平行四边形.

所以 \(D_1M\parallel C_1F\).

因为 \(D_1M\parallel AE\)，所以 \(AE\parallel C_1F\)，于是 \(A,E,F,C_1\) 四点共面，所以点 \(C_1\) 在平面 \(AEF\) 内.
\end{solution}
\begin{problem}
已知函数 \(f(x)=x^3-kx+k^2\).

\begin{enumerate}
    \item 讨论 \(f(x)\) 的单调性；
    \item 若 \(f(x)\) 有三个零点，求 \(k\) 的取值范围.
\end{enumerate}

\end{problem}

\begin{solution}
（1）因为
\(f'(x)=3x^2-k\).

当 \(k\le 0\) 时，\(f'(x)\ge 0\)，所以 \(f(x)\) 在 \(\R\) 上递增.

当 \(k>0\) 时，令 \(f'(x)>0\)，解得
\(x>\sqrt{\frac{k}{3}}\)或\(x<-\sqrt{\frac{k}{3}}\)，
令 \(f'(x)<0\)，解得
\(-\sqrt{\frac{k}{3}}<x<\sqrt{\frac{k}{3}}\).

所以 \(f(x)\) 在 \(\left( -\infty,-\sqrt{\frac{k}{3}} \right)\) 递增，在 \(\left( -\sqrt{\frac{k}{3}},\sqrt{\frac{k}{3}} \right)\) 递减，在 \(\left( \sqrt{\frac{k}{3}},+\infty \right)\) 递增.

综上，\(k\le 0\) 时，\(f(x)\) 在 \(\R\) 上递增；\(k>0\) 时，\(f(x)\) 在 \(\left( -\infty,-\sqrt{\frac{k}{3}} \right)\) 递增，在 \(\left( -\sqrt{\frac{k}{3}},\sqrt{\frac{k}{3}} \right)\) 递减，在 \(\left( \sqrt{\frac{k}{3}},+\infty \right)\) 递增.

（2）由（1）得，当 \(k>0\) 时，\(f(x)\) 的极小值为 \(f\left( \sqrt{\frac{k}{3}} \right)\)，极大值为 \(f\left( -\sqrt{\frac{k}{3}} \right)\).

若 \(f(x)\) 有三个零点，只需
\examdisplaycases{}{
k>0,\\
f\left( \sqrt{\frac{k}{3}} \right)<0,\\
f\left( -\sqrt{\frac{k}{3}} \right)>0.
}

解得
\(0<k<\frac{4}{27}\).

故 \(k\in \left( 0,\frac{4}{27} \right)\).
\end{solution}
\begin{problem}
已知椭圆 \(C:\dfrac{x^2}{25}+\dfrac{y^2}{m^2}=1\)（\(0<m<5\)）的离心率为 \(\dfrac{\sqrt{15}}{4}\)，\(A,B\) 分别为 \(C\) 的左、右顶点.

\begin{enumerate}
    \item 求 \(C\) 的方程；
    \item 若点 \(P\) 在 \(C\) 上，点 \(Q\) 在直线 \(x=6\) 上，且 \(|BP|=|BQ|\)，\(BP\perp BQ\)，求 \(\triangle APQ\) 的面积.
\end{enumerate}

\end{problem}

\begin{solution}
（1）由
\(e=\frac{c}{a}\)
得
\(e^2=1-\frac{b^2}{a^2}\)，
即
\(\frac{15}{16}=1-\frac{m^2}{25}\)，
所以
\(m^2=\frac{25}{16}\).

故椭圆 \(C\) 的方程是
\(\frac{x^2}{25}+\frac{16y^2}{25}=1\).

（2）由（1）知 \(A(-5,0)\)，设 \(P(s,t)\)，\(Q(6,n)\).

根据对称性，只需考虑 \(n>0\) 的情况. 此时 \(-5<s<5\)，\(0<t\le \frac{5}{4}\).

因为 \(|BP|=|BQ|\)，所以
\((s-5)^2+t^2=n^2+1 \circled{1}\)
又因为 \(BP\perp BQ\)，所以
\(s-5+nt=0 \circled{2}\)
且
\(\frac{s^2}{25}+\frac{16t^2}{25}=1 \circled{3}\)

联立 \circled{1}\circled{2}\circled{3}，得
\examdisplayarray{\left\{} {l}{
s=3,\\
t=1,\\
n=2
}{
\right.
\quad \text{或} \quad
\left\{
\begin{array}{l}
s=-3,\\
t=1,\\
n=8.
\end{array}
\right.
}

当 \(s=3\)，\(t=1\)，\(n=2\) 时，\(P(3,1)\)，\(Q(6,2)\)，而 \(A(-5,0)\)，
\(\overrightarrow{AP}=(8,1)\)，\(\overrightarrow{AQ}=(11,2)\).

所以 \(S_{\triangle APQ}=\frac{1}{2}\sqrt{AP^2AQ^2-\left( \overrightarrow{AP}\cdot\overrightarrow{AQ} \right)^2}=\frac{1}{2}\left| 8\times 2-11\times 1 \right|=\frac{5}{2}\).

当 \(s=-3\)，\(t=1\)，\(n=8\) 时，\(P(-3,1)\)，\(Q(6,8)\)，而 \(A(-5,0)\)，
\(\overrightarrow{AP}=(2,1)\)，\(\overrightarrow{AQ}=(11,8)\).

所以
\(S_{\triangle APQ}=\frac{1}{2}\left| 2\times 8-11\times 1 \right|=\frac{5}{2}\).

综上，\(\triangle APQ\) 的面积是 \(\dfrac{5}{2}\).
\end{solution}

\begin{problem}
在直角坐标系 \(xOy\) 中，曲线 \(C\) 的参数方程为 \(\left\{
\begin{array}{l}
x=2-t-t^2,\\
y=2-3t+t^2
\end{array}
\right.\)（\(t\) 为参数且 \(t\neq 1\)），\(C\) 与坐标轴交于 \(A,B\) 两点.

\begin{enumerate}
    \item 求 \(|AB|\)；
    \item 以坐标原点为极点，\(x\) 轴正半轴为极轴建立极坐标系，求直线 \(AB\) 的极坐标方程.
\end{enumerate}

\end{problem}

\begin{solution}
（1）当 \(x=0\) 时，可得 \(t=-2\)（\(1\) 舍去），代入 \(y=2-3t+t^2\)，可得 \(y=12\).

当 \(y=0\) 时，可得 \(t=2\)（\(1\) 舍去），代入 \(x=2-t-t^2\)，可得 \(x=-4\).

所以曲线 \(C\) 与坐标轴的交点为 \((-4,0)\)、\((0,12)\)，故
\(|AB|=\sqrt{(-4)^2+12^2}=4\sqrt{10}\).

（2）由（1）可得直线 \(AB\) 过点 \((0,12)\)、\((-4,0)\)，所以直线 \(AB\) 的方程为
\(\frac{y}{12}-\frac{x}{4}=1\)，
即
\(3x-y+12=0\).

由
\(x=\rho\cos\theta\)，\(y=\rho\sin\theta\)，
可得直线 \(AB\) 的极坐标方程为
\(3\rho\cos\theta-\rho\sin\theta+12=0\).
\end{solution}
\begin{problem}
设 \(a,b,c\in \R\)，\(a+b+c=0\)，\(abc=1\).

\begin{enumerate}
    \item 证明：\(ab+bc+ca<0\)；
    \item 用 \(\max\{a,b,c\}\) 表示 \(a,b,c\) 中的最大值，证明：\(\max\{a,b,c\}\ge \sqrt[3]{4}\).
\end{enumerate}

\end{problem}

\begin{solution}
证明：（1）因为 \(a+b+c=0\)，所以
\((a+b+c)^2=0\)，
即
\(a^2+b^2+c^2+2ab+2ac+2bc=0\).

所以
\(2ab+2ac+2bc=-(a^2+b^2+c^2)\).

又因为 \(abc=1\)，所以 \(a,b,c\) 均不为 \(0\)，从而
\(2ab+2ac+2bc=-(a^2+b^2+c^2)<0\).

故
\(ab+ac+bc<0\).

（2）不妨设 \(a\le b<0<c<\sqrt[3]{4}\)，则
\(ab=\frac{1}{c}>\frac{1}{\sqrt[3]{4}}\).

又因为 \(a+b+c=0\)，所以
\(-a-b=c<\sqrt[3]{4}\).

而
\(-a-b\ge 2\sqrt{ab}>2\sqrt{\frac{1}{\sqrt[3]{4}}}=\sqrt[3]{4}\)，
与假设矛盾.

故
\(\max\{a,b,c\}\ge \sqrt[3]{4}\).
\end{solution}
